%=============================================================================!
% 10.5  STARTING A SIMULATION                                                 !
%=============================================================================!
\section{Starting a simulation}
\label{sec:md-start}

A run needs a starting state: positions inside the cell and a velocity
for every atom. This section builds crystals from their cells, places
them at the spacing their model prefers, and gives the atoms velocities
with no drift, as \cref{sec:ne-atoms} promised.

\subsection*{Positions from a crystal}

A crystal is a cell and the fractional coordinates of the atoms in it.
The face-centred cubic lattice of argon has a cubic cell of side $a$ with
an atom at each corner and at the centre of each face. A corner is shared
by the eight cells that meet there and a face by two, so each cell holds
$8\times\frac18 + 6\times\frac12 = 4$ atoms, at
\begin{equation*}
   (0, 0, 0) , \quad (0, \tfrac12, \tfrac12) , \quad (\tfrac12, 0, \tfrac12) ,
   \quad (\tfrac12, \tfrac12, 0) ,
\end{equation*}
and graphite's hexagonal cell of \cref{sec:md-image} has four, at $(0, 0,
\frac14)$, $(0, 0, \frac34)$, $(\frac13, \frac23, \frac14)$ and $(\frac23,
\frac13, \frac34)$, two in each layer, with the layers shifted against each
other\cite{TrucanoChen1975}. A supercell of $n_1\times n_2\times n_3$ cells
has the lattice vectors multiplied by $n_1$, $n_2$ and $n_3$, and an atom
at $\vb{s}$ in the cell numbered $(k_1, k_2, k_3)$ has the fractional
coordinates $((s_1 + k_1)/n_1, (s_2 + k_2)/n_2, (s_3 + k_3)/n_3)$ in the
supercell. Libraries such as ASE\cite{Larsen2017} build such structures
too, and Notebook~10 checks the positions built here against ASE's.

The spacing matters. A crystal started away from the spacing at which its
model's energy is least is squeezed or stretched, and pushes on its cell;
Chapter~14 measures that push, the pressure. For the Lennard-Jones model
with the energy switched off between $2\sigma$ and $2.5\sigma$
(\cref{sec:pe-cutoffs}), the energy per atom of the face-centred cubic
lattice is least at $a_0 = 1.54916\sigma$, with neighbours $a_0/\sqrt2 =
1.0954\sigma$ apart (from $(0, 0, 0)$ to $(0, \frac12, \frac12)a_0$, a
distance $a_0\sqrt{\frac14 + \frac14}$), and an energy of
$-7.819\varepsilon$ per atom, found by
minimising the sum over pairs with respect to $a$. With $\sigma =
\SI{3.4}{\angstrom}$ that is $a_0 = \SI{5.267}{\angstrom}$, the spacing
used for argon below.

\subsection*{Velocities}

Temperature, which will set how much kinetic energy to give, is defined in
Chapter~12. Until then we choose the kinetic energy directly. The function
\texttt{starting\_velocities} of \texttt{mdlab.md} draws each component at
random, removes the drift of the centre of mass, and scales the velocities
to the chosen kinetic energy:

\begin{code}
    m = np.asarray(masses, dtype=float)
    v = rng.normal(size=(len(m), 3)) / np.sqrt(m)[:, None]
    v -= (m @ v) / m.sum()
    kinetic = 0.5 * mv2_to_energy * float(m @ (v * v).sum(1))
    return v * np.sqrt(kinetic_energy / kinetic)
\end{code}

\noindent Line 2 draws numbers from the normal distribution of width 1,
the bell-shaped spread of \texttt{rng.normal} that Chapter~12 derives, and
divides each atom's by $\sqrt{m_i}$: a component $x/\sqrt{m_i}$ carries
the kinetic energy $\frac12m_i(x/\sqrt{m_i})^2 = \frac12x^2$, whatever the
mass, so light and heavy atoms carry the same kinetic energy on average.
Line 3 subtracts the velocity of the centre of mass, $\vb{P}/M$
(\cref{sec:ne-atoms}). Line 4 is the kinetic energy, with
\texttt{mv2\_to\_energy} the factor \texttt{MV2\_TO\_EV} of
\cref{eq:en-mv2}, or 1 in reduced units, and line 5 scales every velocity by the same factor, which keeps the total
momentum zero. In a periodic cell the angular momentum is not conserved
(\cref{sec:md-periodic}) and is left alone.

\subsection*{Where the energy goes}

Started on its lattice with only kinetic energy, a crystal does not keep
it. \Cref{fig:md-run}(a) of the next section follows 500 argon atoms
started at $a_0$ with \SI{0.02}{\electronvolt} of kinetic energy per atom:
within \SI{0.1}{\pico\second} the kinetic energy falls to half, it
overshoots to \SI{0.0057}{\electronvolt} at \SI{0.14}{\pico\second}, and
between 0.5 and \SI{1}{\pico\second} it averages
\SI{0.0104}{\electronvolt}, as it does over the last \SI{10}{\pico\second}
of a \SI{20}{\pico\second} run. Near its minimum the crystal vibrates in
normal modes (\cref{sec:os-modes}), and a vibration shares its energy
equally, on average, between kinetic and potential (\cref{sec:os-energy}).
Three of the modes move every atom together and have no restoring force,
but the velocities carry no drift, so they hold no energy. All the energy started as kinetic, so on average
half moves into the potential energy. The share left as kinetic energy is
0.518 of the start at \SI{0.02}{\electronvolt} per atom, 0.506 at
\SI{0.005}{\electronvolt} and 0.502 at \SI{0.002}{\electronvolt}: the
excess over a half shrinks with the energy, as the wells look more nearly
parabolic to smaller vibrations. A run started this way therefore needs a
first stretch, discarded before measuring, in which the energy settles
between its two forms.

\begin{keyidea}
Positions come from a crystal's cell and fractional coordinates,
replicated into a supercell at the spacing the model prefers. Velocities
are drawn at random, stripped of the drift of the centre of mass, and
scaled to the chosen kinetic energy; started on its lattice, a crystal
then passes about half of that energy to its potential energy.
\end{keyidea}

\notebook{10}{build a supercell and give it velocities}

\begin{selfcheck}
How many atoms has a supercell of $5\times5\times5$ face-centred cubic
cells, and how long is its side for $a_0 = \SI{5.267}{\angstrom}$?
\end{selfcheck}
